\chapter{Lie Groups and Reflections}
\label{ch:7}
\section{Reflections}
\label{sec:7-1}
Let $u$ be a non-null pure split quaternion. The reflection in the
plane $Q$-orthogonal to $u$ is
\[
r_u(v)=-uvu^{-1}.
\]
It sends $u$ to $-u$ and fixes any pure $v$ orthogonal to $u$, since
then $uv=-vu$. Null vectors cannot be used in this formula because
they are not invertible. Composing two such reflections yields a
conjugation action and hence a determinant-one orthogonal map.

For $H=F+G$ one has $H^2=2I$, $H^{-1}=H/2$, and
\[
HFH^{-1}=G,\quad HGH^{-1}=F,\quad HiH^{-1}=-i.
\]
This exchanges the two spacelike axes and reverses the timelike axis.
It preserves $Q$ and has determinant one on the pure space, but lies
outside the time-orientation-preserving component.
See Section~\ref{sec:7-2} for the determinant issue in the spinor space.
See \cite[Section~7]{KocikFeynman}; notation follows the book conventions.

\section{Krawtchouk matrices as group elements}
\label{sec:7-2}
Set $H=F+G=\begin{pmatrix}1&1\\1&-1\end{pmatrix}$. Then
\[
FH=HG,\qquad GH=HF,\qquad H^2=2I.
\]
These elementary identities generate the Krawtchouk intertwining and
reflection symmetries under symmetric powers.

A qualification to the supplied source is necessary:
$\det H=-2$, and $\det(cH)=-2c^2<0$ for every nonzero real scalar $c$.
Thus no real scalar normalization puts $H$ in $SL(2,\mathbb R)$.
Likewise $\det F=\det G=-1$.
The conjugations are valid in $GL(2,\mathbb R)$ (or its projective
quotient); $F,G$ are also traceless Lie algebra elements. These are
distinct statements, and must not be conflated.
Over $\mathbb C$, $(\mathrm{i}/\sqrt2)H$ has determinant one, where
$\mathrm{i}$ denotes the scalar imaginary unit.

The lift to degree-$N$ polynomials is developed once in
Section~\ref{sec:8-1}; orthogonal normalization is treated in
Section~\ref{sec:14-1}. Neither construction requires the real matrix
$H$ to belong to $SL(2,\mathbb R)$.
See \cite[Section~7]{KocikFeynman}; notation follows the book conventions.

